\documentclass[aspectratio=169,11pt]{beamer}
\usepackage{amsmath,amssymb,graphicx}
\usetheme{default}
\usefonttheme{serif}
\setbeamertemplate{navigation symbols}{}
\setbeamercolor{normal text}{fg=black,bg=white}
\setbeamercolor{structure}{fg=black}
\setbeamercolor{frametitle}{fg=black,bg=white}
\setbeamercolor{title}{fg=black}
\setbeamerfont{frametitle}{size=\Large,series=\bfseries}
\setbeamerfont{title}{size=\LARGE,series=\bfseries}
\setbeamertemplate{itemize item}{\small$\bullet$}
\setbeamertemplate{itemize subitem}{\small$\circ$}
\setbeamertemplate{footline}{\hspace{6mm}\scriptsize Spectral graph theory\hfill\insertframenumber\hspace{6mm}\vspace{3mm}}
\setlength{\parskip}{.45em}
\newcommand{\one}{\mathbf 1}
\newcommand{\E}{\mathbb E}
\newcommand{\F}{\mathbb F}
\newcommand{\Prob}{\mathbb P}
\newcommand{\R}{\mathbb R}
\newcommand{\Afig}[2][]{\includegraphics[#1]{present/slides/assets/#2.pdf}}
\newcommand{\source}[1]{\par\vfill{\scriptsize\color{black!65}#1}}
\title{Spectral graph theory and exact hidden-group recovery}
\subtitle{Examples, a matching-product estimate, and a recovery theorem}
\author{Angelo Farfan \and Carlos Shum \and Ananth Venkatesh}
\institute{Monadic Burrito}
\date{November 15, 2026}
\hypersetup{pdftitle={Spectral graph theory and exact hidden-group recovery},pdfauthor={Angelo Farfan, Carlos Shum, Ananth Venkatesh}}
\begin{document}

\begin{frame}
\titlepage
\end{frame}

\begin{frame}{The route from a graph to a theorem}
\begin{enumerate}
\item Describe a graph by a matrix.
\item Read structure from its eigenvectors and eigenvalues.
\item Understand why random matching entries cancel after centering.
\item Recover exact hidden groups from their neighbor-count equations.
\end{enumerate}
The background assumes only basic algebra. Each matrix operation will be
illustrated on a small graph.
\source{Background reading: Pan, Nie and Wang, \href{https://math.mit.edu/research/highschool/primes/circle/documents/2025/Brian-Nathan-Daniel-Final.pdf}{\emph{An Introduction to Spectral Graph Theory}} (2025). Figures and worked examples here are newly drawn.}
\end{frame}

\section{Graphs and matrices}
\begin{frame}{Vertices, edges, and degree}
A \emph{simple undirected graph} $G=(V,E)$ consists of vertices and
unordered pairs of distinct vertices called edges.

The \emph{degree} $d(v)$ is the number of neighbors of $v$.
\begin{center}\Afig[width=.48\textwidth]{path-three}\end{center}
In this path, the degrees are $1,2,1$.
\[
 \sum_{v\in V}d(v)=2|E|=4.
\]
Each edge contributes once at each endpoint. This proves the
\emph{handshaking identity} for every finite graph.
\end{frame}

\begin{frame}{Four graph families to recognize}
\begin{center}\Afig[width=\textwidth]{graph-families}\end{center}
\begin{itemize}
\item $P_5$: two endpoints of degree one; three vertices of degree two.
\item $C_5$: every vertex has degree two.
\item $K_4$: every vertex has degree three.
\item $K_{1,4}$: one vertex has degree four; four have degree one.
\end{itemize}
A graph is \emph{$d$-regular} if every vertex has degree $d$.
\end{frame}

\begin{frame}{Walks, paths, and connected components}
A \emph{walk} is a sequence of vertices joined by edges.
Its length is the number of edges traversed.
\begin{center}\Afig[width=.47\textwidth]{path-three}\end{center}
\begin{itemize}
\item $1,2,3$ is a path of length two: no vertex repeats.
\item $1,2,1$ is a walk of length two, but not a path.
\item A graph is \emph{connected} if every pair can be joined by a walk.
\end{itemize}
A \emph{connected component} is a maximal connected part of the graph.
Matrices can count walks and even count components.
\end{frame}

\begin{frame}{The adjacency matrix}
Number the vertices $1,\ldots,n$. Define
\[
 A_{ij}=\begin{cases}1,&\{i,j\}\in E,\\0,&\text{otherwise}.\end{cases}
\]
\begin{columns}[c]
\column{.45\textwidth}\Afig[width=\linewidth]{path-three}
\column{.50\textwidth}
\[
 A(P_3)=\begin{pmatrix}0&1&0\\1&0&1\\0&1&0\end{pmatrix}.
\]
\end{columns}
The matrix is symmetric. Its diagonal is zero.
The sum of row $i$ is the degree of vertex $i$.
\end{frame}

\begin{frame}{Matrix multiplication adds the neighbors' values}
Put a number $x_i$ at each vertex. Then
\[
 (Ax)_i=\sum_{j:\{i,j\}\in E}x_j.
\]
For the path on three vertices,
\[
 \begin{pmatrix}0&1&0\\1&0&1\\0&1&0\end{pmatrix}
 \begin{pmatrix}2\\5\\7\end{pmatrix}
 =\begin{pmatrix}5\\9\\5\end{pmatrix}.
\]
The middle vertex adds $2+7=9$.
The endpoints each see only the middle value, $5$.

This local addition rule is the starting point for both spectral methods
and the exact recovery algorithm later in the talk.
\end{frame}

\begin{frame}{Powers of the adjacency matrix count walks}
The entry $(A^2)_{ij}$ counts walks of length two from $i$ to $j$:
\[
 (A^2)_{ij}=\sum_k A_{ik}A_{kj}.
\]
For $P_3$,
\[
 A^2=\begin{pmatrix}1&0&1\\0&2&0\\1&0&1\end{pmatrix}.
\]
There are two length-two returns from vertex $2$:
$2,1,2$ and $2,3,2$.

More generally, $(A^\ell)_{ij}$ counts length-$\ell$ walks. Induction
works because the last step chooses the final neighbor.
\end{frame}

\begin{frame}{An eigenvector reproduces itself under neighbor addition}
A nonzero vector $x$ is an \emph{eigenvector} with \emph{eigenvalue}
$\lambda$ if $Ax=\lambda x$.
\[
 A(P_3)\begin{pmatrix}1\\\sqrt2\\1\end{pmatrix}
 =\begin{pmatrix}\sqrt2\\2\\\sqrt2\end{pmatrix}
 =\sqrt2\begin{pmatrix}1\\\sqrt2\\1\end{pmatrix}.
\]
The three eigenvalues of $A(P_3)$ are $\sqrt2,0,-\sqrt2$.
For example, $(1,0,-1)^{\mathsf T}$ has eigenvalue zero:
the two contributions at the middle vertex cancel.

For a real symmetric matrix, all eigenvalues are real and an orthonormal
eigenbasis exists: the spectral theorem.
\end{frame}

\begin{frame}{Three visible modes on a cycle}
\begin{center}\Afig[width=.96\textwidth]{cycle-modes}\end{center}
The displayed numbers are vector coordinates, not vertex names.
At every vertex, add the two neighboring numbers.

For $C_4$, the full adjacency spectrum is $2,0,0,-2$.
The second zero mode is $(0,1,0,-1)^{\mathsf T}$.
\end{frame}

\begin{frame}{Regularity produces an exact eigenvector}
If $G$ is $d$-regular, then
\[
 A\one=d\one.
\]
Every row adds exactly $d$ ones.

If $G$ is also bipartite, give one side value $+1$ and the other $-1$.
Every neighbor lies on the opposite side, so
\[
 Ax=-dx.
\]
For example, $C_4$ is $2$-regular and bipartite, explaining its eigenvalues
$2$ and $-2$ without a determinant calculation.

Other eigenvalues describe variations away from these simple modes.
\end{frame}

\section{The Laplacian and partitioning}
\begin{frame}{The graph Laplacian measures local disagreement}
Let $D=\operatorname{diag}(d(1),\ldots,d(n))$. The Laplacian is
\[
 L=D-A,\qquad (Lx)_i=\sum_{j\sim i}(x_i-x_j).
\]
For $P_3$,
\[
 L=\begin{pmatrix}1&-1&0\\-1&2&-1\\0&-1&1\end{pmatrix},\quad
 L\begin{pmatrix}1\\0\\-1\end{pmatrix}
 =\begin{pmatrix}1\\0\\-1\end{pmatrix}.
\]
The constant vector has eigenvalue zero. The full Laplacian spectrum is
$0,1,3$.
\end{frame}

\begin{frame}{The basic energy identity}
For every real vector $x$,
\[
 x^{\mathsf T}Lx=\sum_{\{i,j\}\in E}(x_i-x_j)^2\ge0.
\]
\textbf{Proof.} Each edge contributes $x_i^2+x_j^2$ to
$x^{\mathsf T}Dx$ and $2x_ix_j$ to $x^{\mathsf T}Ax$.
Subtract and sum.

For $P_3$ and $x=(1,0,-1)^{\mathsf T}$,
\[
 x^{\mathsf T}Lx=(1-0)^2+(0-(-1))^2=2.
\]
A small energy means that neighboring vertices usually receive similar
values.
\end{frame}

\begin{frame}{Zero eigenvalues count connected components}
By the energy identity,
\[
 Lx=0\quad\Longleftrightarrow\quad
 x\text{ is constant on every connected component}.
\]
One direction follows because every squared difference must vanish.
The other follows directly from $(Lx)_i=\sum_{j\sim i}(x_i-x_j)$.

Thus
\[
 \dim\ker L=\text{number of connected components}.
\]
For a connected graph, zero is a simple Laplacian eigenvalue.
This is an exact algebraic statement about connectivity.
\end{frame}

\begin{frame}{The second eigenvector is the smoothest nonconstant signal}
For a connected graph, write $0=\lambda_1<\lambda_2\le\cdots\le\lambda_n$.
The spectral theorem gives
\[
 \lambda_2=\min_{x\ne0,\;x\perp\one}
 \frac{\sum_{\{i,j\}\in E}(x_i-x_j)^2}{\sum_i x_i^2}.
\]
A minimizing vector is a \emph{Fiedler vector}.
The orthogonality condition excludes constant signals.

Spectral partitioning orders vertices by this vector, then tests threshold
cuts. A slowly varying signal can expose a weak connection between groups.
The best cut and a balanced cut are not automatically the same thing.
\end{frame}

\begin{frame}{Example: two complete graphs joined by one edge}
\begin{center}\Afig[width=.91\textwidth]{spectral-cut}\end{center}
Here the Fiedler vector has one sign on each $K_4$.
Thresholding at zero finds the single bridge.

For this example, $\lambda_2=3-\sqrt7\approx0.354$.
This illustrates the method; it is not a recovery guarantee for every graph.
\end{frame}

\begin{frame}{A cut gives an upper bound on the second eigenvalue}
In the preceding eight-vertex graph, put $x_i=+1$ on the left $K_4$
and $x_i=-1$ on the right. Then $x\perp\one$.

Only the bridge contributes to the energy:
\[
 x^{\mathsf T}Lx=4,\qquad x^{\mathsf T}x=8,
 \qquad\lambda_2\le\frac48=\frac12.
\]
The actual minimizing vector gives the sharper value $3-\sqrt7$.

This calculation explains why a sparse connection can produce a small
Laplacian eigenvalue.
\end{frame}

\begin{frame}{Keep adjacency and Laplacian eigenvalues distinct}
For a $d$-regular graph, $L=dI-A$. Thus
\[
 Ax=\theta x\quad\Longrightarrow\quad Lx=(d-\theta)x.
\]
For irregular graphs, $D$ is not a scalar matrix, so this conversion
is unavailable.

A \emph{spectral gap} means a specified separation between eigenvalues;
one must say which matrix and which eigenvalues.

The recovery theorem at the end does not require a planted adjacency
eigenvalue to be isolated from the other adjacency eigenvalues.
\end{frame}

\section{Random lifts and centered matching products}
\begin{frame}{A random lift replaces each vertex by a fiber}
Replace each vertex of a fixed base graph by $N$ copies.
Replace each base edge by an independent uniform perfect matching
between its two fibers.
\begin{center}\Afig[width=.69\textwidth]{lift-example}\end{center}
A lift of a $d$-regular base remains $d$-regular.
The picture is one example, not an asymptotic experiment.
\end{frame}

\begin{frame}{Old modes and new modes}
A vector constant on each fiber can be lifted from the base graph.
Neighbor addition on such vectors is exactly the base adjacency action.

Therefore every base eigenvalue is inherited by every lift.
These are the \emph{old} eigenvalues; the orthogonal complement carries
the \emph{new} spectrum.

Bounding the new spectrum requires control of long random walks.
Their products involve entries of the same random matchings, so the
entries cannot be treated as independent.
\source{Bordenave (2020), random-lift spectral theorem and Proposition 28.}
\end{frame}

\begin{frame}{A perfect matching is a permutation matrix}
A uniform permutation $\sigma$ of $\{1,\ldots,N\}$ defines
\[
 M_{ij}=\mathbf 1_{\{\sigma(i)=j\}},\qquad \E M_{ij}=N^{-1}.
\]
Every row and every column has exactly one entry equal to one.

For $r$ prescribed matches in distinct rows and columns,
\[
 \Prob(\text{all occur})=\frac{(N-r)!}{N!}=\frac1{(N)_r},
 \quad (N)_r=N(N-1)\cdots(N-r+1).
\]
Two prescribed matches in the same row cannot both occur.
\end{frame}

\begin{frame}{Centering: a two-cell calculation}
\begin{center}\Afig[width=.65\textwidth]{requested-cells}\end{center}
\[
 \E[(M_{11}-N^{-1})(M_{22}-N^{-1})]
 =\frac1{N(N-1)}-\frac1{N^2}=\frac1{N^2(N-1)}.
\]
In the same row, the expectation is instead $-N^{-2}$.
The pattern of requested cells determines how much cancellation occurs.
\end{frame}

\begin{frame}{Compatible conditioning leaves a smaller uniform matrix}
After fixing $r$ compatible matches, delete their rows and columns.
The remaining matrix $U$ is a uniform permutation matrix of size $N-r$.
\begin{center}\Afig[width=.50\textwidth]{depleted-matrix}\end{center}
Keep the original centering: $Z_{ij}=NU_{ij}-1$.
Then each remaining row satisfies
\[
 \sum_v Z_{iv}=N-(N-r)=r.
\]
\end{frame}

\begin{frame}{A row sum gives a positive recurrence}
Let $\mu_j=\E\prod_{i=1}^j Z_{ii}$, with $\mu_0=1$ and $r+j\le N$.
Multiplying a row sum by $\prod_{i<j}Z_{ii}$ gives
\[
 (N-r-j+1)\mu_j
 =(r+2j-2)\mu_{j-1}+(j-1)\mu_{j-2}.
\]
Omit the final term when $j=1$.

The identity for a used column is
$Z_{ii}Z_{ji}=-Z_{ii}-Z_{ji}-1$.
All recurrence coefficients on the right are nonnegative.
This retains cancellation before any absolute-value estimate.
\source{The report proves the depleted recurrence. At $r=0$, the corresponding recurrence is prior work of Collins--Nagatsu.}
\end{frame}

\begin{frame}{A numerical check of the recurrence}
Take $N=6$ and condition on $r=2$ compatible matches.
The remaining matrix has size four.
\[
 \mu_1=\frac{2}{4}=\frac12,\qquad
 3\mu_2=4\mu_1+\mu_0=3,\qquad \mu_2=1.
\]
Directly,
\[
 \E[(6U_{11}-1)(6U_{22}-1)]
 =\frac{36}{4\cdot3}-\frac{12}{4}+1=1.
\]
The recurrence is positive even though individual centered entries can
have either sign.
\end{frame}

\begin{frame}{The matching-product estimate: what the gain means}
Consider $a$ distinct requested cells, with arbitrary centered and raw
powers and independent permutation colors.
Let $h$ be the centered singletons whose row and column contain no other
requested cell of that color; let $b$ count the other centered singletons.

For $2\le a\le\sqrt N$, the report proves
\[
 |\E P|\le\sqrt e\,2^bN^{-a}
 \left(\frac{\sqrt{a-1}}{\sqrt N-\sqrt{a-1}}\right)^h.
\]
Each isolated singleton contributes an extra small factor of order
$\sqrt{a/N}$. Repeated powers introduce no extra factor of two.
\source{The report's scalar mixed bound implies Bordenave's Proposition 28 form. Exact attained envelopes are also proved there.}
\end{frame}

\begin{frame}{The biregular model uses half-edge matchings}
\begin{center}\Afig[width=.68\textwidth]{biregular-matching}\end{center}
For $n_L$ left vertices of degree $d_L$ and $n_R$ right vertices of
degree $d_R$, both half-edge sets have size
$H=n_Ld_L=n_Rd_R$.

The random cross block is a uniform permutation matrix on these $H$
half-edges. Apply the product estimate with $N=H$ and centering $1/H$.
\source{Brito--Dumitriu--Harris (2022): published Lemma 5.6, preprint Lemma 7.}
\end{frame}

\begin{frame}{How the estimate enters the spectral argument}
A \emph{non-backtracking walk} never immediately reverses its previous
edge. Its expansion leads to products of matching entries.
\begin{enumerate}
\item Expand a high matrix norm moment into walk patterns.
\item Bound the expectation of each matching product.
\item Count patterns and control tangles.
\item Relate non-backtracking roots to adjacency eigenvalues.
\end{enumerate}
For any matrix $C$, including a nonnormal one,
\[
 \|C\|^{2s}\le\operatorname{tr}((CC^*)^s).
\]
\source{The report proves the product input. The graph path enumeration and tangle bounds remain cited inputs; it does not infer a new graph spectral theorem from the product bound alone.}
\end{frame}

\section{Exact hidden-group recovery}
\begin{frame}{A new question: can exact groups be recovered?}
Suppose each vertex has exactly two neighbors in its own group and
one in the other group. The labels of the groups have been erased.
\begin{center}\Afig[height=2.9cm,width=.64\textwidth,keepaspectratio]{two-groups}\end{center}
This is an exact neighbor-count rule, not an average preference.
Can the graph alone reveal the hidden grouping?
\end{frame}

\begin{frame}{An equitable partition satisfies an exact matrix equation}
A partition is \emph{equitable} if all vertices of group $i$ have the
same number $Q_{ij}$ of neighbors in group $j$.

Let $J_{vi}=1$ when vertex $v$ belongs to group $i$, and zero otherwise.
Then
\[
 AJ=JQ.
\]
For two groups with internal degree $a$ and cross degree $b$,
\[
 Q=\begin{pmatrix}a&b\\b&a\end{pmatrix},\qquad
 A\tau=(a-b)\tau,
\]
where $\tau$ is $+1$ on one group and $-1$ on the other.
The hidden groups create an exact algebraic signal.
\end{frame}

\begin{frame}{The random model and its sufficient condition}
Fix upper bounds $K$ on the number of groups and $D$ on vertex degree.
The true count table $Q=(d_{ij})$ is fixed, has connected type support,
and has fixed positive size proportions obeying
$N_i d_{ij}=N_j d_{ji}$: cross-edge counts agree from both ends.

Each internal regular graph and each cross biregular graph is sampled
independently and uniformly among simple graphs with those counts.

The theorem requires weights satisfying
\[
 q_{ij}=q_{ji},\qquad 0\le q_{ij}<\tfrac12,\qquad
 \sum_jd_{ij}q_{ij}=1.
\]
The weights certify applicability; the algorithm is not given them.
If all degrees equal $d\ge3$, use $q_{ij}=1/d$.
\end{frame}

\begin{frame}{The main recovery theorem, in words}
Under the preceding model and weight condition, as the graph grows:
\begin{itemize}
\item With probability tending to one, every equitable partition into at
most $K$ nonempty groups only merges whole planted groups.
\item The planted partition is the unique finest such partition.
\item From the graph and $K,D$, the algorithm lists every valid partition;
selecting the one with most groups recovers the plant.
\end{itemize}
Neither the true table, the number of groups, nor their sizes is supplied.
Class names are arbitrary, so recovery is up to renaming.
\source{The report's recovery extension states the precise theorem and supplies the proof. This is a fixed-parameter, noiseless random-model guarantee.}
\end{frame}

\begin{frame}{A concrete consequence: every fixed regular base}
Take any fixed connected simple $d$-regular base graph, $d\ge3$.
An independent random $N$-lift satisfies the weight condition.
\begin{center}\Afig[width=.63\textwidth]{lift-example}\end{center}
With probability tending to one, the finest partition with at most $K$
groups reveals the fibers, provided $K$ is at least the base order.
The algorithm need not be told which base generated the graph.
\end{frame}

\begin{frame}{Why the bound on the number of groups is essential}
If every vertex forms its own group, the partition is always equitable:
each neighbor count is simply an adjacency entry.

So ``find the finest equitable partition'' would be trivial if the
number of groups were unrestricted.

The theorem asks for the finest partition with at most a supplied
\emph{fixed} number $K$ of nonempty groups.

On the structural event, every other admissible answer is obtained by
merging whole planted groups, so the largest class count identifies the
plant unambiguously.
\end{frame}

\begin{frame}{The algorithm first tries finitely many count tables}
With $K,D$ fixed, there are only finitely many nonnegative integer tables
of order at most $K$ and row sums at most $D$.
\begin{enumerate}
\item Try each candidate table $R$.
\item Solve a modular version of $AY=YR$.
\item Enumerate the solutions represented by a complete kernel basis.
\item Keep only one-hot, nonempty partitions passing the integer equation.
\end{enumerate}
Finally, retain a verified partition with the largest number of groups.

The number of tables may be large. It is part of the fixed-parameter
constant, not a claim that the method is immediately practical.
\end{frame}

\begin{frame}{Finite-field arithmetic organizes the search}
Arithmetic modulo three works on the clock $0,1,2$:
\[
 2+1=0,\qquad 2\cdot2=1\quad\text{in }\F_3.
\]
For each proposed $R$, solve the linear system
\[
 AY-YR=0\quad\text{over a fixed finite field}.
\]
Every genuine partition survives reduction modulo that prime.
Some modular solutions are false candidates.

Therefore every output is checked again using the original \emph{integer}
neighbor counts. Modular coincidence is never sufficient for acceptance.
\end{frame}

\begin{frame}{A complete kernel basis is a certificate of coverage}
If a finite-field solution space has dimension $h$, a basis describes
all its $p^h$ elements using $h$ coefficients in $\F_p$.

The algorithm certifies the \emph{complete} right kernel, even when the
matrix is nonsymmetric. It then enumerates its candidates and verifies
their integer neighbor counts.

Finding one attractive solution would not prove that all valid partitions
had been considered. Completeness is what justifies selecting a finest one.

The randomized computation is always correct; randomness affects its
running time. Such an algorithm is called \emph{Las Vegas}.
\end{frame}

\begin{frame}{Why the expected search cost is nearly quadratic}
The key probabilistic estimate is a \emph{raw moment} bound:
\[
 \sup_n\E\big[p^{h_R(A)}\big]<\infty
\]
for every fixed candidate matrix $R$ and field, where
$h_R(A)=\dim\{Y:AY=YR\}$ over that field.

It counts the average number of modular candidates, including rare
inputs with very large solution spaces.

Combined with the certified sparse computation, it gives expected
$\widetilde O(n^2)$ bit work and expected $\widetilde O(n)$ working bits.
The tilde hides logarithmic factors; fixed-parameter constants may be large.
\end{frame}

\begin{frame}{The proof has three logically separate parts}
\begin{enumerate}
\item \textbf{Identifiability.} Entropy controls large alternative
partitions; a sparse-support count excludes small alternatives.
\item \textbf{Search size.} Finite-field channel moments bound the actual
expected number of modular candidates, including exceptional inputs.
\item \textbf{Certified computation.} A complete-kernel algorithm and
integer checks find exactly the valid partitions at the stated cost.
\end{enumerate}
Correct listing holds on every input. Identifying the particular hidden
plant is a separate event whose probability tends to one in the model.
\end{frame}

\begin{frame}{What ``without a spectral gap'' means here}
For the two-group example, $A\tau=(a-b)\tau$ holds exactly.
The algorithm uses the neighbor-count equations themselves.

It does not require the eigenvalue $a-b$ to be separated numerically
from the remaining eigenvalues, nor a one-dimensional real eigenspace.

This does not establish noise tolerance. Changing edges can destroy the
exact equations.

Likewise expected nearly quadratic work is not a quadratic worst-case
bound. Some inputs can have many valid answers and a long search.
\end{frame}

\begin{frame}{A sharp ambiguity example from the full research manuscript}
For equal groups of size $N$, internal degree $2$ and cross degree $1$,
choose the hidden balanced grouping uniformly.
\begin{center}\Afig[width=.59\textwidth]{ambiguity-rates}\end{center}
The ambiguity probability is $\sim4/N^2$; the optimal average error is
$\sim2/N^2$. These are asymptotic formulas, not finite-size error bars.
\source{Proved in the full 60-page research manuscript, two-class ambiguity section. This additional result is distinct from the report's condensed core recovery proof.}
\end{frame}

\begin{frame}{Boundary cases show why the assumptions matter}
\textbf{Degree two.} An unconditioned lift of a known cycle is a collection
of cycles. Components admit different shifts and orientations of their
hidden labels. Optimal hidden-fiber success tends to zero.

\textbf{Unequal degrees.} A $K_{3,12}$ lift has exponentially many compatible
fiber partitions, even after the three left fibers are revealed.
Under a uniform remaining hidden prior, optimal success is exponentially small.

Thus minimum degree three by itself does not suffice for arbitrary
nonregular templates. The positive theorem uses the full weight condition.
\source{These separate information bounds are proved in the full research manuscript. Connected conditioning changes the cycle experiment.}
\end{frame}

\begin{frame}{The contribution and its predecessors}
The proposed advance is a proved recovery and expected-cost guarantee
across the specified fixed-parameter class, with unknown tables and sizes.

The algebraic recovery idea has direct predecessors:
\begin{itemize}
\item Fellows--Koblitz (1993): hidden fibers in $K_4$ covers, Kid Krypto.
\item Bernstein--Cottaar--Lange (2023): Gaussian elimination and grouping
equal affine expressions to recover those fibers.
\end{itemize}
The manuscript credits these algorithms and classical counting, entropy,
finite-field moments, and certified linear algebra.
Historical firstness is not certified.
\end{frame}

\begin{frame}{What the theorem does and does not promise}
\begin{itemize}
\item \textbf{Exact random model:} fixed group and degree bounds;
independent uniform simple factors; exact neighbor counts.
\item \textbf{Statistical recovery:} succeeds with probability tending to
one as the model grows; correct listing is unconditional.
\item \textbf{Work bound:} expected nearly quadratic bit work;
no benchmark-speed or small-constant guarantee.
\item \textbf{Limits:} no noise theorem, no growing-parameter uniformity,
no general clustering guarantee for arbitrary networks.
\end{itemize}
The work has internal mathematical reviews and finite checks.
These do not constitute formal verification or external peer review.
\end{frame}

\begin{frame}{References: background and the matching estimate}
\small
\begin{itemize}
\item B.~Pan, N.~Nie, D.~Wang.
\emph{An Introduction to Spectral Graph Theory}. MIT PRIMES Circle, 2025.
\item C.~Bordenave. \emph{A new proof of Friedman's second eigenvalue theorem
and its extension to random lifts}. Ann. Sci. ENS 53 (2020), 1393--1439.
\item G.~Brito, I.~Dumitriu, K.~D.~Harris.
\emph{Spectral gap in random bipartite biregular graphs and applications}.
CPC 31 (2022), 229--267.
\item B.~Collins, M.~Nagatsu. \emph{Weingarten calculus for centered random
permutation matrices}. Algebraic Combinatorics 9 (2026), 465--479.
\end{itemize}
\end{frame}

\begin{frame}{References and proof locations: recovery}
\small
\begin{itemize}
\item M.~Fellows, N.~Koblitz. \emph{Kid Krypto}. CRYPTO '92 proceedings,
published 1993.
\item D.~J.~Bernstein, J.~Cottaar, T.~Lange. IPCC evaluation, Chapter 3
of \emph{Report on evaluation of KpqC candidates}, 2023, Section 3.1.
\item \emph{Spectral-gap report}: original matching-product proofs followed
by the condensed recovery extension and its technical proofs.
\item \emph{Unknown-frame recovery in sparse block-regular graphs},
November 15, 2026: complete source paper, additional ambiguity,
old/new-spectrum, and information-obstruction results.
\end{itemize}
The full source paper and this talk distinguish proved statements,
cited inputs, internal checks, and historical-priority limits.
\end{frame}
\end{document}
